\section{Hyperbolic Space: The World That Keeps Opening Up}

We need to live inside hyperbolic space for a while, because the whole crystal sits in it. The problem is that our eyes and intuitions were built for flat space, so hyperbolic space feels alien at first. There is a famous picture that makes it visible, and once you have it, everything gets easier.

\subsection{The Poincare disk}

You cannot draw the whole of hyperbolic space on a flat page honestly, because there is too much of it. But there is a clever, faithful way to squeeze all of it into a finite circle. It is called the Poincare disk, and it is the device behind those Escher prints where angels and devils shrink as they march toward the edge of a circle.

Here is how to read it. The entire infinite hyperbolic plane is shown inside one round disk. The catch is the scale. Near the center, distances look normal. As you move toward the rim, everything gets drawn smaller and smaller. A shape near the edge looks tiny on the page, but in the real hyperbolic world it is exactly the same size as a shape in the middle. The shrinking is an artifact of cramming infinity into a circle.

Two consequences follow, and they are the whole feel of the place.

First, \textbf{the rim is infinitely far away}. As you walk toward the edge, your steps get drawn shorter and shorter on the page, so you never arrive. There is always more space between you and the rim. The edge is not a wall. It is the horizon at infinity.

Second, \textbf{there is always more room out there}. Because room opens up exponentially, the closer you get to the rim, the more cells fit in each ring. The crowd at the edge is bigger than everything inside it. This is why the Escher figures get so numerous and small near the boundary.

\subsection{Straight lines that look bent}

In hyperbolic space, the shortest path between two points, the ``straight line,'' is drawn in the Poincare disk as an arc that bows gently toward the center. This looks wrong to a flat-space eye, but it is correct: that bowed arc really is the shortest route. A useful image is that the center of the disk is ``downhill'' and the rim is ``uphill,'' so the cheapest path always dips toward the middle a little.

This single fact will matter later, when we send signals across the crystal. The straight, cheap paths go through the deep interior, not along the crowded rim.

\subsection{No center, really}

The disk seems to have a special center, but that is a lie of the drawing. Hyperbolic space looks exactly the same from every point. You can slide the whole picture so that any point becomes the center, and nothing about the geometry changes. There is no privileged spot and, crucially, no privileged direction. Every place is like every other place, and at every place all directions are alike.

That last property, sameness of every direction, is the seed of a deep payoff. Our universe has no preferred direction either. Later we will see that a crystal built in this directionless curved space inherits that directionlessness, which is what saves the whole theory from a fatal problem that wrecks crystals built in flat space.

\subsection{Holding the picture}

Keep this image in your pocket for the rest of the walkthrough: a circle, packed with cells that shrink toward an infinitely distant rim, the same in every direction, with the cheap paths bowing through the middle, and always, always more room out at the edge.

\subsection{Finite, but never finished}

If the rim of hyperbolic space is infinitely far away, is the universe infinite? In this model, no, not at any moment. The crystal is always finite, a definite number of cells, however large. What is endless is its growth: there is no largest size it ever reaches, because it keeps adding cells at the frontier forever. The infinitely distant rim of the disk is the limit the crystal grows toward and never arrives at. So the universe is finite and unbounded at once, finite at every instant, with no final size, growing without end.

\textbf{Takeaway:} the Poincare disk shows all of hyperbolic space inside a circle, where the rim is infinitely far, room explodes toward the edge, straight lines bow inward, and every point and direction is alike. That sameness of direction is a gift we will cash in later.
